\begin{abstract}
Post-training---supervised fine-tuning followed by reinforcement learning from human feedback---turns a raw next-token predictor into a chatbot that follows instructions, refuses harmful requests, and adopts a conversational register. Although the behavioural effects of this transformation are well documented, its representational effects remain poorly understood: does post-training carve new concepts into the model, or does it merely re-weight and re-route the conceptual vocabulary that pre-training already laid down? This question matters for alignment, since interventions targeting ``safety features'' only make sense if such features exist as discrete entities that post-training has put there.

This thesis investigates that question by directly comparing the feature dictionaries of the Gemma 3 1B base and instruction-tuned checkpoints, using Sparse Autoencoders (SAEs) from the Gemma Scope 2 release as the interpretability substrate. At residual-stream layer 13, features extracted from the two models are matched on two complementary criteria---decoder-direction cosine similarity and per-sequence activation correlation---and classified as preserved, modified, PT-exclusive, or IT-exclusive. Selected features are then labelled by two independent LLMs, Claude Haiku 4.5 and GPT-4.1-nano, filtered by inter-model consensus, probed against safety and instruction-following prompts using Welch's $t$-test on differential activations, and finally subjected to causal verification via norm-calibrated activation steering.

The central finding supports what this thesis calls the ``wrapper hypothesis.'' At moderate matching thresholds ($\tau_d = 0.7$, $\tau_a = 0.3$), only about 3\% of features are strictly preserved while about 49\% are modified, meaning that they share a clear geometric counterpart across models but differ in how they fire. Crucially, when the safety- and instruction-relevant IT features surfaced by the targeted analysis are traced back to the matching results, the majority turn out to have a PT match rather than being IT-exclusive. Steering experiments confirm this picture causally: features that the matching pipeline labels as IT-exclusive produce graded behavioural changes when injected into the base model, but so do features that already existed in the base. On this evidence, post-training looks less like a sculptor adding new material and more like a conductor rebalancing an existing orchestra.

The contribution is threefold: first, a reproducible end-to-end SAE comparison pipeline for paired base/IT checkpoints, including an explicit treatment of threshold sensitivity; second, a cross-model labelling protocol that uses inter-LLM agreement as a confidence filter rather than relying on a single annotator; and third, an application of Gemma Scope 2 SAEs to PT-vs-IT comparison on Gemma 3 1B, with results that alignment-by-steering work can use directly. The findings suggest that interpretability work aimed at safety should focus less on hunting for newly created alignment features and more on understanding how post-training reorganises pre-existing ones.
\end{abstract}
